\section{综合练习}

\answer{AA01}
最终向量的总奇偶性必为偶数。对任意满足该约束的 $n$ 维比特向量，其前 $n-1$ 位唯一决定生成过程，因而概率均为 $2^{-(n-1)}$；不满足约束的向量概率为零。

任取 $n-1$ 个坐标并指定其取值，未指定的最后一个坐标由偶数奇偶性唯一决定。因此这个指定事件的概率恰为 $2^{-(n-1)}$，等于 $n-1$ 个公平比特各自概率之积。这证明任意这样的子组都独立。

全部为零的概率也是 $2^{-(n-1)}$，而若 $n$ 个比特独立应为 $2^{-n}$，两者不同，所以整体不独立。奇偶约束只有同时观察全部坐标才显现。

\answer{AA02}
对任意整数 $k$，
\[
\ind_{\{k\equiv r\ (3)\}}
=\frac13\sum_{j=0}^2\omega^{j(k-r)}.
\]
若 $k-r$ 是3的倍数，三项全为1；否则它们是全部三次单位根，和为零。取期望并利用独立性，
\[
\E\omega^S=\left(\frac{1+\omega}{2}\right)^n
 =2^{-n}e^{in\pi/3}.
\]
另一个非平凡项是其乘上相应相位后的共轭，故
\[
\boxed{\Prob(S\equiv r\ (3))
=\frac13+\frac{2^{1-n}}3
\cos\left(\frac{n\pi}{3}-\frac{2\pi r}{3}\right).}
\]
这个表达式也覆盖 $n=0$。对 $n=8,r=0$，余弦为 $-1/2$，答案为 $1/3-1/768=85/256$。直接加上 $S=0,3,6$ 的二项概率得到 $(1+56+28)/256$，一致。

\answer{AA19}
每次交换改变排列奇偶性，因此奇数步不可能复原。把两步视为一轮：连续交换同一对位置的概率为 $1/2$，排列不变；交换不同对的位置产生两个方向的三循环，各以概率 $1/4$ 发生。

从偶排列出发，两步链只在三个偶排列中运动，转移矩阵的对角元为 $1/2$，非对角元为 $1/4$。设第 $m$ 轮后回到初态的概率为 $u_m$，则
\[
u_{m+1}=\frac12u_m+\frac14(1-u_m)
=\frac14+\frac14u_m,\qquad u_0=1.
\]
减去不动点 $1/3$ 得到
\[
\boxed{\Prob(\text{第 }t\text{ 步复原})=
\begin{cases}
0,&t\text{ 为奇数},\\
\dfrac13+\dfrac23\,4^{-m},&t=2m.
\end{cases}}
\]
偶数子序列趋于 $1/3$，奇数子序列恒为零，故不存在统一极限。状态有限且转移对称，并不能自动消除周期性。

\answer{AA22}
对每个候选中心位置引入指示变量。圆周上三个位置互不相同，依次要求中心是 B、两个邻居是 A，概率为
\[
\frac nN\frac m{N-1}\frac{m-1}{N-2}.
\]
有 $N$ 个候选中心，故圆周期望为
\[
\boxed{\frac{nm(m-1)}{(N-1)(N-2)}}.
\]
线性排列只有 $N-2$ 个非端点中心，每个中心仍有相同模式概率，期望为
\[
\boxed{\frac{nm(m-1)}{N(N-1)}}.
\]
若 $m<2$ 或 $n=0$，两式都为零，符合定义。代入 $(m,n)=(4,3)$，分别为 $6/5$ 和 $6/7$。相邻中心的指示变量可能相关，线性期望并不要求它们独立。

\answer{AA03}
横坐标为 $x$ 附近的竖直切片长为 $2\sqrt{R^2-x^2}$。除以总面积，得到
\[
f_X(x)=\frac{2}{\pi R^2}\sqrt{R^2-x^2}\,\ind_{\{|x|<R\}}.
\]
圆盘中央的切片更长，所以投影更集中于0附近。对称性给出 $\E X=0$，且
\[
\E|X|=\frac4{\pi R^2}\int_0^R x\sqrt{R^2-x^2}\,\dd x
=\frac{4R}{3\pi}.
\]
计算二阶矩可用旋转对称：$\E X^2=\E Y^2=\E(X^2+Y^2)/2$。半径平方在 $[0,R^2]$ 上均匀，故 $\E(X^2+Y^2)=R^2/2$，最终
\[
\boxed{\Var X=R^2/4}.
\]
投影密度中的切片权重正是不能把坐标直接视为均匀变量的原因。

\answer{AA07}
以概率1所有点互不相同。沿圆周读取类型，能够严格分开当且仅当 A 点组成一个连续的圆周块，B 点组成其补块。充分性来自连接两个类型分界间隙的弦：适当选择弦端点使其避开样本点，弦所在直线分开两块。必要性来自直线与圆周最多有两个交点，两侧各对应一个圆弧；若类型交替出现超过两次，就不可能分开。

固定一个编号圆周起点后，类型排列均匀分布于 $\binom{m+n}{m}$ 个选择。连续 A 块可从 $m+n$ 个位置开始；由于两类都非空，这些选择互不相同。因此
\[
\boxed{\Prob(\text{可分})=\frac{m+n}{\binom{m+n}{m}}}.
\]
$m=n=3$ 时为 $6/20=3/10$。概率只由循环类型次序决定，实际角间距不进入计数。

\answer{AA09}
标准二维正态的方向均匀，条件于正象限后极角 $\Theta$ 在 $(0,\pi/2)$ 上均匀，且 $Z=\cot\Theta$。对 $z>0$，
\[
F_Z(z)=\Prob(\Theta\ge\arctan(1/z))
=\frac2\pi\arctan z.
\]
对 $z\le0$，分布函数为0。求导得到
\[
f_Z(z)=\frac2{\pi(1+z^2)}\quad(z>0).
\]
$F_Z(1)=1/2$，所以中位数为1；但
\[
\int_0^\infty z f_Z(z)\,\dd z
=\frac2\pi\int_0^\infty\frac z{1+z^2}\,\dd z=\infty.
\]
分子分母都具有全部有限矩，并不保证比值有有限均值；分母接近零时会产生重尾。

\answer{AA10}
条件于 $S=s>0$ 后，方向仍均匀。交换对称性给出 $\E[X^2\mid S=s]=s/3$。记
\[
a=\E[X^4\mid S=s],\qquad b=\E[X^2Y^2\mid S=s].
\]
符号对称使含奇次幂的混合矩为零；旋转后的 $(X+Y)/\sqrt2$ 与 $X$ 同分布，故
\[
a=\frac{2a+6b}{4},\qquad a=3b.
\]
再展开恒等式 $s^2=(X^2+Y^2+Z^2)^2$，取条件期望得 $s^2=3a+6b=15b$。因此
\[
\E[X^2Y^2\mid S]=\frac{S^2}{15},\quad
\Var(X^2\mid S)=\frac{4S^2}{45},\quad
\Cov(X^2,Y^2\mid S)=-\frac{2S^2}{45}.
\]
这些等式几乎处处成立。固定总平方和后，一个坐标的能量增大必须由其余坐标补偿，解释了条件负相关；它并不与未条件化时的独立性矛盾。

\answer{AA13}
把时刻排序为 $U_1<U_2<U_3$，几乎必然严格不等。最短距离为 $\min(U_2-U_1,U_3-U_2)$。对 $0\le d\le L/2$，事件 $D>d$ 等价于两个内部间隔都大于 $d$。作平移
\[
v_1=U_1,\qquad v_2=U_2-d,\qquad v_3=U_3-2d.
\]
可行域变成 $0<v_1<v_2<v_3<L-2d$，体积为 $(L-2d)^3/6$；排序密度为 $6/L^3$。故
\[
F_D(d)=
\begin{cases}
0,&d<0,\\
1-(1-2d/L)^3,&0\le d\le L/2,\\
1,&d>L/2.
\end{cases}
\]
由尾积分，
\[
\E D=\int_0^{L/2}(1-2d/L)^3\,\dd d=\frac L8.
\]
平移扣除的是两个受约束的内部间隔，而非四个含端点的间隔，这是尺度因子2的来源。

\answer{AA16}
以概率1最小值由唯一的下标 $J$ 取得。固定 $J=j$，令其余剩余量为 $r_i=X_i-U>0$。对应联合密度为
\[
\lambda^n\exp\!\left[-\lambda\!\left(nu+\sum_{i\ne j}r_i\right)\right]
=\frac1n\,(n\lambda e^{-n\lambda u})
 \prod_{i\ne j}\lambda e^{-\lambda r_i}.
\]
这是离散下标 $J$ 的均匀概率、$U$ 的指数密度以及 $n-1$ 个独立指数剩余量密度的乘积。因此 $U\sim\operatorname{Exp}(n\lambda)$，且剩余量与 $U$ 独立。极差是这些剩余量的最大值，于是
\[
\Prob(D\le d)=
\begin{cases}0,&d<0,\\(1-e^{-\lambda d})^{n-1},&d\ge0.\end{cases}
\]
由尾积分并令 $y=1-e^{-\lambda d}$，
\[
\E D=\frac1\lambda\int_0^1\frac{1-y^{n-1}}{1-y}\,\dd y
=\frac1\lambda\sum_{k=1}^{n-1}\frac1k.
\]
令 $H_{n-1}=\sum_{j=1}^{n-1}1/j$，则 $\E[V\mid U]=U+H_{n-1}/\lambda$。这里独立性来自指数密度的精确分解，不能只由零协方差推断。

\answer{AA17}
当 $|\rho|<1$ 时，可写 $Y=\rho X+\sqrt{1-\rho^2}\,Z$，其中 $Z$ 与 $X$ 独立标准正态。条件于 $X$，交叉项的均值为零，所以
\[
\E[Y^2\mid X]=\rho^2X^2+1-\rho^2.
\]
由 $\E X^2=1$、$\E X^4=3$，
\[
\E[X^2Y^2]=\rho^2\E X^4+(1-\rho^2)\E X^2
=1+2\rho^2.
\]
因此 $\Cov(X^2,Y^2)=2\rho^2$，两者的方差都为 $3-1=2$，得到
\[
\boxed{\Corr(X^2,Y^2)=\rho^2}.
\]
边界 $|\rho|=1$ 时 $Y=\rho X$ 几乎必然，上式直接成立。平方使 $\rho$ 与 $-\rho$ 给出相同结果，所以相关符号丢失。非线性条件均值仍可准确计算，并不等于把线性回归系数简单平方后忽略常数项。

\answer{AA20}
先设 $r>0$。沿 $a$ 的方向取极轴，$U=\cos\Theta$ 均匀分布于 $[-1,1]$，且
\[
D^2=R^2+r^2-2RrU.
\]
故 $D$ 的密度为 $d/(2Rr)$，支持 $|R-r|<d<R+r$。于是
\[
\E\frac1D=\frac{R+r-|R-r|}{2Rr}
=\boxed{\frac1{\max(R,r)}}.
\]
$r=0$ 时 $D=R$，结论继续成立。即使固定点落在球面上，零距离只发生于零概率事件，倒数的一阶积分仍有限。

若 $r\ne R$ 且 $r>0$，
\[
\E\frac1{D^2}
=\frac1{2Rr}\log\frac{R+r}{|R-r|}.
\]
$r=R$ 时下限为零，积分变成 $\int_0^{2R}\dd d/d$，发散。一个几何奇点可以不影响一阶可积性，却破坏二阶可积性。

\answer{AA06}
任意新两次试验相同的概率为 $p^2+q^2>0$，故等待时间有几何尾。记末次为成功、失败时的剩余均值分别为 $e_S,e_F$。首步分解为
\[
e_S=1+qe_F,\qquad e_F=1+pe_S.
\]
解得 $e_S=(1+q)/(1-pq)$、$e_F=(1+p)/(1-pq)$。开始时先做一次试验，因此
\[
\boxed{\E T=1+pe_S+qe_F=\frac{2+pq}{1-pq}}.
\]
再令 $a_S,a_F$ 为两状态下最终由SS吸收的概率。递推 $a_S=p+qa_F$、$a_F=pa_S$，得到
\[
\boxed{\Prob(\text{SS停止})
=pa_S+qa_F=\frac{p^2(1+q)}{1-pq}}.
\]
把 $p,q$ 对换得到FF概率，二者相加为1。公平情形均值为3，两类吸收各占一半。停止前最后一次结果决定哪种下一结果会吸收，因此不能只用一个“尚未结束”状态。

\answer{AA08}
记第 $i$ 次结果为 $X_i$，令 $I_i=\ind_{\{X_i=H,X_{i+1}=T\}}$，$1\le i\le n-1$，并令 $a=p(1-p)$。则 $N=\sum I_i$，各指标均值为 $a$、方差为 $a(1-a)$。

相隔至少两个起点的指标依赖于不相交的抛掷，所以独立；相邻起点不可能同时出现HT，因为共享字符不能同时为H和T。因此
\[
\Cov(I_i,I_{i+1})=-a^2.
\]
恰有 $n-2$ 对相邻指标，得到
\[
\boxed{\E N=(n-1)a,\qquad
\Var N=(n-1)a-(3n-5)a^2}.
\]
代入 $n=10,a=1/4$，均值为 $9/4$，方差为 $11/16$。把 $N$ 当作二项变量会忽略重叠约束，并高估方差。

\answer{AA12}
均匀分布在对称环游下保持不变，故每时刻都有 $\E Y_t=0$、$\E Y_t^2=1/2$。令 $\lambda=\cos(2\pi/n)\in(-1,1)$。余弦加法公式给出
\[
\E[Y_{t+1}\mid X_t]=\lambda Y_t.
\]
迭代条件期望得 $\E[Y_k\mid X_0]=\lambda^kY_0$，所以
\[
\boxed{\Cov(Y_0,Y_k)=\tfrac12\lambda^k}.
\]
由平稳性逐项展开方差，
\[
\boxed{\Var S_T=\frac12\left[T+
2\sum_{k=1}^{T-1}(T-k)\lambda^k\right]}.
\]
除以 $T$ 后，因 $\sum_{k\ge1}|\lambda|^k<\infty$，可对 $(1-k/T)\ind_{\{k<T\}}\lambda^k$ 使用绝对可和控制，得到
\[
\lim_{T\to\infty}\frac{\Var S_T}{T}
=\frac12+\frac{\lambda}{1-\lambda}
=\frac{1+\lambda}{2(1-\lambda)}.
\]
例如 $n=6,T=4$ 时 $\lambda=1/2$，有限和为 $33/8$，并非独立观测给出的2。

\answer{AA04}
阶乘中的因子2不少于因子5，因此尾零数为
\[
z(n)=\sum_{j\ge1}\left\lfloor\frac n{5^j}\right\rfloor.
\]
它单调不减，且 $z(5K)\ge K$。在整数区间 $[0,5K]$ 二分寻找第一个满足 $z(n)\ge K$ 的 $n_*$。若 $z(n_*)>K$，则目标被一次跳跃越过，无解；若相等，返回该常值区间。

对每个 $m\ge0$，$5m,\ldots,5m+4$ 中所有 $\lfloor n/5^j\rfloor$ 都不变，而从 $5m+4$ 到 $5m+5$ 至少新增一个因子5。因此每个可达值恰对应一个长度5的整数块，算法成功时返回 $n_*,\ldots,n_*+4$。$K=0$ 时得到0至4，同样成立。

一次计算 $z(n)$ 只需反复除以5，使用 $O(\log(n+2))$ 次整数运算。二分总计 $O(\log^2(K+2))$ 次整数运算。$z(100)=20+4=24$，所以 $K=24$ 的解为100至104。$z(120)=28$，$z(125)=25+5+1=31$，故29不可达到。

\answer{AA05}
利用
\[
\binom nk=\frac{(n-k+1)\cdots n}{1\cdots k},
\]
可先用对称性把 $k$ 换成 $\min(k,n-k)$。逐个处理分子、分母的因子：剥去所有2和5，分别累计分子减分母的净指数 $a,b$；其余部分都与10互素，分子单位相乘、分母单位乘其模10逆元，得到 $u\in\{1,3,7,9\}$。这些逆元分别为 $1,7,3,9$，可直接验证。

二项式系数是整数，所以净指数 $a,b\ge0$。令 $z=\min(a,b)$，则末个非零数字为
\[
\boxed{u\,2^{a-z}5^{b-z}\pmod{10}}.
\]
这是在局部剥离后使用单位群中的合法逆元；原分母含2或5时并没有模10逆元，不能直接相除。各指数可用循环剥因子求得，模幂可用周期或快速幂。$k=0$ 时空乘积给出1。

例如 $\binom{20}{7}=77520$，净指数 $a=4,b=1$，单位部分 $u=9$。删去一个10后，$9\cdot2^3\equiv2\pmod{10}$，故答案为2。这里列出完整数值只用于核验，并非算法所需。

\answer{AA11}
位移 $(a,b)$ 的线段含有 $\gcd(|a|,|b|)-1$ 个内部格点。理由是将 $(a,b)$ 除以最大公因数后得到本原整数方向，线段上的整数点只能按该基本步长排列。因此可见当且仅当最大公因数为1。

水平方向只允许相邻点，共 $M(N+1)$ 对；竖直方向有 $N(M+1)$ 对。对 $a,b\ge1$ 的绝对位移，每种斜率方向有 $(M+1-a)(N+1-b)$ 对，两种方向都要计入。可见无序对总数为
\[
V=M(N+1)+N(M+1)
+2\sum_{\substack{1\le a\le M,\ 1\le b\le N\\\gcd(a,b)=1}}
(M+1-a)(N+1-b).
\]
所求概率为 $V/\binom{(M+1)(N+1)}2$。在 $M=N=2$ 时，轴向共有12对；合法正位移为 $(1,1),(1,2),(2,1)$，斜向贡献16对。故 $V=28$，概率为 $28/36=7/9$。

\answer{AA14}
先计不限制大小顺序的正整数三元组，总数为 $\binom{N-1}{2}$。若第三坐标违背严格三角不等式，则 $z\ge x+y$。令 $s=x+y$，有 $2\le s\le\lfloor N/2\rfloor$，每个 $s$ 有 $s-1$ 对正整数 $(x,y)$，因此此类个数为 $\binom{\lfloor N/2\rfloor}{2}$。

由于坐标均为正，不可能同时有两个坐标分别不小于其余两者之和，三类坏事件互斥。有序三角形数故为
\[
A(N)=\binom{N-1}{2}-3\binom{\lfloor N/2\rfloor}{2}
=\frac{N^2}{8}+O(N).
\]
三个边长互异时，每个无序解对应6个有序解。出现相等坐标的解至多 $O(N)$ 个，每个的轨道大小可能为3或1。因此
\[
T(N)=\frac{A(N)}6+O(N)=\frac{N^2}{48}+O(N).
\]
在 $N=12$ 时，$z<6$ 且 $z\ge4$。$z=4$ 给 $(4,4,4)$；$z=5$ 给 $(2,5,5)$、$(3,4,5)$，共3个。它们分别产生1、3、6个排列，共10个有序解，说明简单除以6不能精确计数。

\answer{AA15}
共轭数 $\beta=(3-\sqrt5)/2$ 满足 $0<\beta<1/2$，且 $\alpha,\beta$ 都是 $z^2-3z+1=0$ 的根。令
\[
L_0=2,\quad L_1=3,\quad L_n=3L_{n-1}-L_{n-2}.
\]
归纳或根公式给出 $L_n=\alpha^n+\beta^n\in\mathbb Z$。因此 $L_n-\alpha^n=\beta^n\in(0,1/2)$，它就是唯一最近整数，且误差精确等于 $\beta^n$。同时 $\lfloor\alpha^n\rfloor=L_n-1$。

模2递推为 $L_n\equiv L_{n-1}+L_{n-2}$，从 $(L_0,L_1)\equiv(0,1)$ 出发得到周期 $(0,1,1)$，故 $L_n$ 当且仅当 $3\mid n$ 时为偶数。

误差随 $n$ 严格递减。解 $\beta^n<10^{-6}$ 得 $n>6\log 10/(-\log\beta)$，右端约为14.35，所以最小为15。也可直接比较 $\beta^{14}\approx1.4072\times10^{-6}$ 与 $\beta^{15}\approx5.3749\times10^{-7}$。

\answer{AA21}
积非零，不能含0；因子5、7不能出现，否则会整除72。单个数位不可能等于72。两位数中乘积72仅能用8与9，数位和17，不满足任一目标。

若有三位，将数位从小到大排列。乘积72的全部三数位多重集为
\[
(1,8,9),\ (2,4,9),\ (2,6,6),\ (3,3,8),\ (3,4,6).
\]
可按最小数位依次为1、2、3分类验证完整性。最小数位为4时，余下两位均不小于4，却须乘积等于18，不可能；最小数位至少5时，乘积至少为125，也不可能。

数位和14只保留 $(2,6,6)$ 与 $(3,3,8)$，其最小排列分别266、338，所以答案为266。数位和18只保留 $(1,8,9)$，答案为189。任何四位及以上正整数都比三位答案大。删去189中的1虽然不改变积，却把数位和18改成17，破坏约束。

\answer{AA18}
若 $\log2$ 是代数数，它又不为零，所给定理会推出 $e^{\log2}=2$ 超越，与2为整数矛盾。因此 $\log2$ 超越；同理 $\log3$ 与 $\log4$ 也超越。

假设 $\log2/\log3=a/b$ 是有理数，其中 $a,b$ 为正整数。交叉相乘并取指数得到 $2^b=3^a$。左边只含素因子2，右边只含素因子3，与整数素因子分解的唯一性矛盾，所以该比值无理。

但两个超越数的比值可以有理，例如 $\log4/\log2=2$。上面的证明只断言 $\log2/\log3$ 无理，并没有从所给定理直接推出它超越。不同类型的结论需要各自的论据，不能把“无理”与“超越”互相替换。
